% pxviewer manuscript draft.
% Target genre: the "introduce a GUI application" software paper, modelled on
%   Emsley & Cowtan 2004 (Coot), Pettersen et al. 2021 (ChimeraX),
%   Croll 2018 (ISOLDE).
% For Acta Crystallographica D submission, switch to the IUCr class:
%   \documentclass{iucr}  (iucr.cls from journals.iucr.org) — the section
%   structure below maps onto it directly. For a preprint this compiles as-is.
\documentclass[11pt]{article}
\usepackage[margin=1in]{geometry}
\usepackage{graphicx}
\usepackage{booktabs}
\usepackage{hyperref}
\usepackage[numbers,sort&compress]{natbib}
\usepackage{xcolor}
\newcommand{\pxviewer}{\emph{pxviewer}}
\newcommand{\angstrom}{\AA}
\newcommand{\todo}[1]{\textcolor{red}{[TODO: #1]}}
\setlength{\parskip}{0.4em}
\setlength{\parindent}{0pt}

\title{\pxviewer: an interactive environment for macromolecular model
inspection, validation and refinement}

\author{Christopher Schlicksup%
  \thanks{\todo{affiliation and contact email}}}

\date{}

\begin{document}
\maketitle

\begin{abstract}
\pxviewer{} is an interactive application for inspecting, validating and
refining macromolecular models against experimental data. Atomic models are
read with the Computational Crystallography Toolbox (cctbx) into its
\emph{mmtbx} model framework, so every view, table and selection in the
application is derived from the same validated objects used by the Phenix
suite; structures are rendered by the Mol* toolkit, with Python and the
renderer connected by a live coordinate bridge that streams conformation
changes into the view without reloading topology. The desktop application
organizes inspection around the crystallographer's workflow rather than the
implementation's: residue, atom and geometry-restraint tables act as literal
selections that drive navigation, focus and clipping; MolProbity-style
validation is presented the same way, so a flagged residue is one click from
the camera. Interactive workflows built on the bridge include real-space
refinement into cryo-EM density, X-ray refinement judged against a live
mFo--DFc difference map, ligand fitting and restraint editing. Guided
tutorials walk new users through each of these tasks inside the running
program. \pxviewer{} is distributed as self-contained installers for macOS,
Windows and Linux that require no terminal or package manager.
\end{abstract}

\textbf{Keywords:} molecular graphics; model building; structure validation;
cryo-EM; crystallographic refinement; interactive refinement.

%% =====================================================================
\section{Introduction}

Molecular graphics has occupied a special place in protein crystallography
since \emph{FRODO} \citep{frodo1978} and then \emph{O} \citep{o1991} showed
what interactive display could do for model building; the tradition continued
with \emph{Coot} \citep{coot2004, coot2010}, which remains the standard
interactive model-building tool for X-ray crystallography. Along a parallel
track, UCSF \emph{Chimera} and \emph{ChimeraX} \citep{chimera2004,
chimerax2018, chimerax2021} built the field's most capable general-purpose
visualization environment, and \emph{ISOLDE} \citep{isolde2018} demonstrated
that physically realistic interactive simulation could make low-resolution
remodelling both safer and faster.

Two changes have reshaped the technical landscape since these tools were
designed. First, the data side of macromolecular crystallography has
consolidated around cctbx \citep{cctbx2002, phenix2010}, whose
\emph{mmtbx} model hierarchy, geometry-restraint machinery and validation
metrics encode decades of crystallographic correctness; a viewer that reads
models through cctbx inherits all of it. Second, the best molecular renderer
available today runs not on the desktop but in the browser: Mol*
\citep{molstar2021} delivers publication-quality representations, per-atom
color theming and very large structures entirely in WebGL. Between them lies
a gap: cctbx owns the crystallography but renders nothing interactively,
while Mol* renders superbly but understands no crystallography.

\pxviewer{} bridges that gap. It keeps every crystallographic decision on the
Python side --- model I/O through cctbx's \texttt{DataManager}, geometry
restraints from cctbx's own proxies, MolProbity-derived validation --- and
lets Mol* do what it does best. The two halves are joined by a live
coordinate bridge: a persistent channel over which Python streams the model
topology once, then conformation frames, selections, representations,
measurements and camera state as the model and the user's intent evolve.
Because the bridge is symmetric, interactive tools such as refinement and
ligand fitting present their progress in the renderer as it happens rather
than at the end.

A second design commitment shapes the interface. Crystallographers think in
\emph{tables of things} --- residues, atoms, restraints, validation flags ---
and in \emph{selections} of those things. In \pxviewer{} a table row \emph{is}
a selection: clicking a restraint row selects the atoms it restrains, and the
same selection machinery governs viewport picking, camera focus, clipping,
representations and filtering. Guided tutorials use the same state, so the
program can watch what the user does and acknowledge each step only when it
is truly done.

The result is intended to sit naturally beside the tools it learns from:
Coot-like in its crystallographic substance, ChimeraX-like in its rendering
and tutorial-driven onboarding, and ISOLDE-like in its insistence that
refinement be something you watch, not something you wait for.

%% =====================================================================
\section{Overview and architecture}

\begin{figure}[h]
\centering
\fbox{\parbox{0.85\textwidth}{\centering\vspace{2em}
\todo{Figure 1: architecture diagram --- cctbx data layer $\rightarrow$ live
bridge $\rightarrow$ Mol* renderer, with the Qt desktop shell around both.}
\vspace{2em}}}
\caption{\pxviewer{} architecture. All crystallographic logic lives in Python
(cctbx model I/O, geometry restraints, validation, refinement); all rendering
lives in Mol*. The live bridge streams the topology once and then coordinates,
selections, representations, measurements and camera state in both
directions.}
\label{fig:arch}
\end{figure}

The application has three components (Fig.~\ref{fig:arch}).

\textbf{The data layer} is the \pxviewer{} Python package. Atomic models are
read by cctbx's \texttt{DataManager} into an \texttt{mmtbx.model.manager};
the viewer never parses structural data itself. Maps arrive through the same
path for CCP4, MRC and related formats, and reflection data through cctbx's
\emph{miller} arrays --- so experimental quantities (amplitudes, phases,
R-free flags) retain their meaning rather than being flattened into
coordinates. Structures fetched by accession code from the Protein Data Bank
or EMDB take the identical path.

\textbf{The live coordinate bridge} is a versioned WebSocket protocol. On
connect, Python sends the model topology once, as BinaryCIF; afterwards the
wire carries binary coordinate frames (one \texttt{float32} triple per atom),
a compact binary channel for per-atom attributes (e.g.\ B-factor or
local-resolution values used for coloring), and small JSON messages for
selections, representations, measurement primitives, non-covalent
interaction notation, clashes, camera state and background. Atom identity is
positional and fixed for a session: index $i$ on the wire is always atom
$i$ of the topology, which is what makes streamed conformation updates safe.

\textbf{The renderer} is a Mol*-based frontend (TypeScript/React). It applies
each frame to the streamed topology in place --- representations update
coordinates rather than rebuild --- and reports viewport picks back over the
bridge. Representation state (type, coloring, opacity, scope) is owned by
Python via MolViewSpec and replayed to viewers that connect later.

\textbf{The desktop shell} embeds the same frontend in QtWebEngine
(PySide6) beside a native controls window. Users who prefer it can instead
point any browser at a session served by \texttt{pxviewer serve}; the wire
protocol is identical either way.

%% =====================================================================
\section{Inspecting models and data}

\subsection{Loading and viewing structures}

The \emph{Open} action accepts local model files (PDB, mmCIF and other
cctbx-supported formats), density maps and reflection files, or an accession
code to fetch from the PDB/EMDB. Several models can be loaded at once and
independently shown, hidden or made active. Representations follow
MolViewSpec's vocabulary --- cartoon/ribbon, ball-and-stick, spacefill,
surface, carbohydrate --- with coloring by element, chain, secondary
structure, residue properties, arbitrary uniform colors, or any per-atom
attribute (B-factor, occupancy, or values supplied or computed by the user)
mapped through a named palette and domain.

\subsection{Tables are selections}

The Geometry tab presents the model as the crystallographer already thinks
of it: \emph{Components} lists residue groups (residues, ligands, waters);
\emph{Atoms} is a virtualized table of per-atom attributes that remains
responsive past $10^5$ atoms; and one virtualized table per restraint type
lists bonds, angles, dihedrals, chiralities and planarities, read directly
from the cctbx restraint proxies. Every row is a literal selection of the
atoms it describes, and selections are shared state across the whole
application: a row selected in a table highlights those atoms in the
viewport, an atom picked in the viewport selects its row, camera focus and
the clip sphere can be bound to the running selection, and keyboard stepping
walks the table one residue or one atom at a time. Pick granularity is
explicit --- atom-level by default, promoted to residue where the displayed
representation offers no atom-sized target --- and is visible and
overridable in the Selection pane.

\subsection{Measurements, interactions and clipping}

Distance, angle, dihedral and label measurements are drawn between atoms or
selection centroids and track the structure as it moves. Non-covalent
interaction notation (hydrogen bonds, salt bridges, hydrophobic and other
contact kinds) can be supplied explicitly from Python or inferred by Mol*;
steric clashes are computed from the coordinates themselves and drawn in a
distinct ``bad contact'' style. Clipping is unified across slab, sphere and
selection-driven modes, so isolating a binding site is a property of the
selection rather than a separate mechanism.

%% =====================================================================
\section{Validation}

Validation in \pxviewer{} is not a report but an interface. The
\emph{Validation} tab runs the MolProbity-derived checks available through
cctbx --- rotamer outliers, Ramachandran outliers, CaBLAM backbone outliers,
C$\beta$ deviations, cis- and twisted-peptide flags --- each producing a
table of flagged residues. Tables behave exactly like the geometry tables:
a row is a selection of the atoms involved, so the path from ``this residue
is an outlier'' to ``show me'' is one click.

Two further tools consolidate the picture. \emph{Hotspots} aggregates the
per-residue checks into a single per-atom severity field mapped onto the
structure, because a region worth rebuilding usually trips several checks
at once and a picture finds such regions faster than five tables. For cryo-EM
models, \emph{Q-score} \citep{qscore2020} reports per-atom and per-residue
map-model correlation, colorable directly onto the structure. Steric clashes
computed from the coordinates complete the validation set.

%% =====================================================================
\section{Interactive workflows}

The bridge changes what refinement looks like: instead of a batch job whose
result is inspected afterwards, refinement becomes a sequence the user
watches and can interrupt, adjust and resume. Four worked workflows ship in
the application and in its guided tutorials.

\subsection{Real-space refinement into cryo-EM density}

Real-space refinement \citep{rsr2018} minimizes the model against the
density map itself. In \pxviewer{} the cctbx machinery runs on the Python
side while the bridge streams each accepted conformation into the viewport,
so the user watches the fit converge in the map rather than waiting for a
final structure. Restraints, weighting and secondary-structure constraints
are the same ones the Phenix suite applies.

\subsection{X-ray refinement with a live difference map}

X-ray refinement is judged against data. \pxviewer{} computes the
mFo--DFc difference map from the loaded reflections as the model changes and
streams it live: positive and negative contours update around the region
being rebuilt, so the map functions as real-time feedback rather than a
post-hoc verdict. This mirrors the working loop experienced
crystallographers run between refinement rounds, compressed into the
interactive session.

\subsection{Ligand fitting and restraint editing}

The ligand-fitting workflow locates unmodeled difference density and fits a
flexible ligand into it. Restraint edits --- custom bonds, angles or
metal-coordination restraints the monomer library cannot know --- are
authored from the atoms themselves inside the GUI and can be shared as a
PHIL file for use in Phenix refinement, so an interactive session produces
artifacts that survive it.

\subsection{Local resolution}

A deposited cryo-EM resolution is a global average; \pxviewer{} computes
local resolution and colors the structure or map by it, making the
soft-periphery/sharp-core distinction directly visible.

%% =====================================================================
\section{Guided tutorials}

Every workflow above is also a guided tutorial, built into the application's
\emph{Tutorials} menu. A tutorial is a non-modal ``coach'' that loads its
own example data (so its instructions always describe the structure actually
on screen), then steps the user through the task one instruction at a time.
Two properties distinguish it from documentation: each step carries a
predicate over live application state, so the step is acknowledged only when
the user has genuinely done the thing; and a \emph{Show me where} button
flashes the exact widget the step refers to. The tutorials double as
regression tests --- the same step sequences run headlessly in the test
suite --- so they cannot drift out of sync with the interface.

%% =====================================================================
\section{Scripting and reproducibility}

The same bridge that drives the GUI is a public API. A
\texttt{LiveSession} can be created from a model file, an existing
\emph{mmtbx} manager or synthetic coordinates; it exposes selections (by
index, by the Phenix atom-selection language, or by mask), representations,
per-atom coloring, measurements, interaction notation, clashes, camera
state and mouse selection reported back to Python. The desktop app embeds an
IPython console bound to the active model's session, so the GUI and scripts
drive the same object.

Views are reproducible: \emph{Save view} serializes the session --- loaded
objects, representations, colors, clipping, selections, visibility,
background and camera --- to a file that restores the scene exactly, with
optional bundling of the source data for archival. This is intended to make
published figures reproducible artifacts rather than screenshots of a state
no one can reconstruct.

%% =====================================================================
\section{Distribution and availability}

\pxviewer{} is distributed as self-contained installers for macOS
(Apple~Silicon), Windows~10/11 (x86\_64) and Linux (x86\_64), downloadable
from the project releases page. Each packages a complete conda environment
--- Python, cctbx, the chem\_data monomer and validation libraries, Qt and
the built Mol* frontend --- so installation needs neither a terminal nor a
package manager, and the Linux bundle works fully offline. A source
install (a single script that bootstraps conda, builds the environment and
the frontend) supports development on macOS and Linux, and a
\texttt{noarch} conda package is available for managed deployments.

\pxviewer{} is source-available software (see the license accompanying the
distribution); source code, issue tracking and releases are at
\url{https://github.com/cschlick/pxviewer}.

%% =====================================================================
\section{Implementation notes}

\textbf{Atom identity.} Frames, selections and per-atom attributes are
positional throughout: atom $i$ is row $i$ of the topology. The atom set may
not change mid-session, which turns the streaming problem into a fixed-size
contract rather than a reconciliation problem.

\textbf{Packaging.} The conda recipe is the single source of dependency
truth; platform build scripts pack a fresh environment into the artifact
(dmg, NSIS archive-in-installer, tar.xz), build the validation caches inside
it, and trim the chem\_data datasets the application cannot reach (BLAST
databases for eLBOW, the standalone secondary-structure regularization
library). Prefix fixup is deferred to launch/install time so the bundles are
relocatable.

\textbf{Testing.} A headless suite covers the Python layer and GUI
invariants; visual passes run the renderer on all three platforms, including
the tutorial step sequences --- the walkthroughs ship only if they still
work.

%% =====================================================================
\section{Discussion and outlook}

\pxviewer{} deliberately occupies the space its predecessors left between
them. From Coot it takes the premise that crystallographic primitives ---
restraints, monomer chemistry, validation --- should drive the interface;
from ChimeraX it takes the standard for rendering quality, onboarding and
distribution; from ISOLDE it takes the conviction that refinement should be
visible and interruptible. Its distinguishing choices are the ones those
tools could not have made at their inception: cctbx as the sole authority
for crystallographic meaning, a browser-grade renderer as the sole
authority for pixels, and a persistent bidirectional bridge in between.

Natural directions follow from that architecture: multiple simultaneous
viewers on one session, collaboration on shared selections and views,
deeper integration of the live bridge with Phenix refinement programs, and
extension of the tutorial framework into user-authorable walkthroughs for
teaching.

%% =====================================================================
\section*{Acknowledgements}
\todo{acknowledgements --- testers, beta users, funding.}

%% =====================================================================
\begin{thebibliography}{99}

\bibitem{frodo1978} Jones TA (1978). A graphics model building and refinement
system for macromolecules. \emph{J. Appl. Cryst.} \textbf{11}, 268--272.

\bibitem{o1991} Jones TA, Zou J-Y, Cowan SW, Kjeldgaard M (1991). Improved
methods for building protein models in electron density maps and the location
of errors in these models. \emph{Acta Cryst. A} \textbf{47}, 110--119.

\bibitem{coot2004} Emsley P, Cowtan K (2004). Coot: model-building tools for
molecular graphics. \emph{Acta Cryst. D} \textbf{60}, 2126--2132.

\bibitem{coot2010} Emsley P, Lohkamp B, Scott WG, Cowtan K (2010). Features
and development of Coot. \emph{Acta Cryst. D} \textbf{66}, 486--501.

\bibitem{chimera2004} Pettersen EF, Goddard TD, Huang CC, Couch GS,
Greenblatt DM, Meng EC, Ferrin TE (2004). UCSF Chimera --- a visualization
system for exploratory research and analysis. \emph{J. Comput. Chem.}
\textbf{25}, 1605--1612.

\bibitem{chimerax2018} Goddard TD, Huang CC, Meng EC, Pettersen EF, Couch GS,
Morris JH, Ferrin TE (2018). UCSF ChimeraX: meeting modern challenges in
visualization and analysis. \emph{Protein Sci.} \textbf{27}, 14--25.

\bibitem{chimerax2021} Pettersen EF, Goddard TD, Huang CC, Meng EC, Couch GS,
Croll TI, Morris JH, Ferrin TE (2021). UCSF ChimeraX: structure visualization
for researchers, educators, and developers. \emph{Protein Sci.} \textbf{30},
70--82.

\bibitem{isolde2018} Croll TI (2018). ISOLDE: a physically realistic
environment for model building into low-resolution electron-density maps.
\emph{Acta Cryst. D} \textbf{74}, 519--530.

\bibitem{cctbx2002} Grosse-Kunstleve RW, Sauter NK, Moriarty NW, Adams PD
(2002). The Computational Crystallography Toolbox: crystallographic
algorithms in a reusable software framework. \emph{J. Appl. Cryst.}
\textbf{35}, 126--136.

\bibitem{phenix2010} Adams PD, Afonine PV, Bunk\'oczi G, \emph{et al.} (2010).
PHENIX: a comprehensive Python-based system for macromolecular structure
solution. \emph{Acta Cryst. D} \textbf{66}, 213--221.

\bibitem{molstar2021} Sehnal D, Bittrich S, Deshpande M, \emph{et al.} (2021).
Mol* Viewer: modern web app for 3D visualization and analysis of large
biomolecular structures. \emph{Nucleic Acids Res.} \textbf{49}, W431--W437.

\bibitem{molprobity2010} Chen VB, Arendall WB, Headd JJ, \emph{et al.} (2010).
MolProbity: all-atom structure validation for macromolecular crystallography.
\emph{Acta Cryst. D} \textbf{66}, 12--21.

\bibitem{molprobity2018} Williams CJ, Headd JJ, Moriarty NW, \emph{et al.}
(2018). MolProbity: more and better reference data for improved all-atom
structure validation. \emph{Protein Sci.} \textbf{27}, 293--315.

\bibitem{rsr2018} Afonine PV, Poon BK, Read RJ, \emph{et al.} (2018).
Real-space refinement in PHENIX for cryo-EM and crystallography.
\emph{Acta Cryst. D} \textbf{74}, 531--544.

\bibitem{qscore2020} Pintilie G, Zhang K, Su Z, Li S, Schmid MF, Chiu W
(2020). Measurement of atom resolvability in cryo-EM maps with Q-scores.
\emph{Nat. Methods} \textbf{17}, 328--334.

\bibitem{mvs} \todo{MolViewSpec citation if one exists; otherwise cite the
project URL.}

\end{thebibliography}
\end{document}
